\documentclass[journal]{IEEEtran}
\usepackage[T1]{fontenc}
\usepackage[utf8]{inputenc}
\usepackage[spanish]{babel}
\usepackage{csquotes}

% *** CITATION PACKAGES ***
\usepackage[style=ieee,backend=biber]{biblatex}
\addbibresource{bibliografia.bib}

% *** MATH PACKAGES ***
\usepackage{amsmath}

% *** PDF, URL AND HYPERLINK PACKAGES ***
\usepackage{url}
\hyphenation{op-tical net-works semi-conduc-tor}
\usepackage{graphicx}
\usepackage{tikz}
\usetikzlibrary{shapes.geometric, arrows.meta, positioning, calc}
\usepackage{float}
\usepackage{array}
\newcolumntype{L}[1]{>{\raggedright\arraybackslash}p{#1}}
\usepackage{hyperref}

\begin{document}

% paper title
\title{Integración de Inteligencia Artificial en el Análisis Forense Digital}

% author names 
\author{Angie Camila Gutiérrez Trujillo, %
\and Juan Esteban Triviño,  %
\and Juan Miguel Delgado %
}

% The report headers
\markboth{Integración de Inteligencia Artificial en el Análisis Forense Digital}%
{IA y Forensia Digital}

% make the title area
\maketitle
\thispagestyle{empty}
\pagestyle{empty}

%===========================================================
% SECCIÓN 3: RESUMEN
%===========================================================
\begin{abstract}
\normalfont\itshape
El crecimiento exponencial de la información digital y la sofisticación de los vectores de ataque han incrementado significativamente la complejidad de las investigaciones forenses modernas \cite{PadilhaRafael2021AIAe}. Este documento presenta una formulación analítica sobre la integración de la Inteligencia Artificial (IA) en el ecosistema de análisis forense digital. Acotando el alcance temático, se explora cómo el aprendizaje automático y la visión computacional optimizan el triaje de dispositivos móviles \cite{VASILARAS2024301737} y la detección de contenido sintético o \textit{deepfakes} \cite{BharatiNitu2026EddA, Al-ImranMd2025DACF}. Asimismo, se contrastan los fundamentos inmutables de la cadena de custodia frente a la necesidad crítica de implementar modelos de IA Explicable (XAI) \cite{HermosillaPamela2025EAfF, HallStuartW.2022APoC}. Este enfoque resulta imperativo para mitigar la opacidad de los algoritmos de ``caja negra'', garantizando la integridad, la trazabilidad y el valor probatorio de los hallazgos en escenarios judiciales bajo el control constante del factor humano \cite{FAHNDRICH2023301617, 10.1145/3560107.3560114}.
\end{abstract}

\begin{IEEEkeywords}
Forense Digital, Inteligencia Artificial (IA), Explicabilidad (XAI), Deepfakes, Admisibilidad Legal, Internet de las Cosas (IoT).
\end{IEEEkeywords}

%===========================================================
% SECCIÓN 5: INTRODUCCIÓN
%===========================================================
\section{Introducción}
\IEEEPARstart{L}{a} ciencia forense digital, históricamente enfocada en la recuperación metódica y preservación de evidencias estáticas, enfrenta hoy un desafío sin precedentes. El volumen masivo de datos generados por dispositivos inteligentes e infraestructuras críticas \cite{KimSungbum2022Addf}, sumado a la rápida evolución de la IA generativa, han transformado drásticamente tanto la naturaleza de la evidencia como las tácticas empleadas para ocultarla o falsificarla \cite{10.1145/3733102.3733138}. Ante este panorama, la evolución desde métodos puramente manuales hacia el uso de sistemas automatizados e inteligencia artificial resulta imperativa para agilizar la resolución de incidentes y mitigar el problema del \textit{Big Data} forense \cite{ShailjaSingh2026TEoF, QuickDarren2014Ioiv}.

De acuerdo con los alcances temáticos delineados para esta investigación, el presente trabajo se circunscribe a tres campos de aplicación específicos: el triaje y priorización automatizada en dispositivos móviles, la detección de imágenes manipuladas (contenido sintético y \textit{deepfakes}), y la aplicación de Inteligencia Artificial Explicable (XAI) a las conclusiones forenses.

Sin embargo, delegar la correlación e interpretación de evidencia a modelos computacionales introduce riesgos técnicos y legales profundos. Los modelos predictivos de ``caja negra'', que operan con altos niveles de opacidad, dificultan la trazabilidad del razonamiento y la interpretación forense \cite{SharmaDevanshu2026LaEI}. Esta falta de transparencia hace urgente el desarrollo de una IA que no contamine la evidencia original y mantenga al perito humano como validador final \cite{VASILARAS2024301737}. El presente trabajo investiga los vectores de integración de estas tecnologías disruptivas, garantizando la salvaguarda de los principios fundamentales de la investigación forense y su validación legal ante autoridades judiciales.

%===========================================================
% SECCIÓN 6: FORMULACIÓN DEL PROBLEMA
%===========================================================
\section{Formulación del problema}
El procesamiento de evidencia digital basado exclusivamente en la revisión manual o en reglas estáticas resulta técnica y operativamente insuficiente frente a ataques cada vez más sofisticados, como la alteración dinámica de marcas de tiempo \cite{10.1145/3712716.3712727} y la proliferación masiva de contenidos sintéticos indetectables a simple vista \cite{NevesMiguelJoseDas2025SCFD}. Si bien la IA promete revolucionar la automatización del triaje y la detección de anomalías \cite{SerhalCezar2021Mlba}, su naturaleza inherentemente probabilística y la posible introducción de sesgos de entrenamiento riñen con la exactitud determinista exigida en los estrados judiciales \cite{MaccaroneLeeT.2025IACi}. Por lo tanto, el problema central de esta investigación se formula de la siguiente manera:

¿Cómo puede integrarse la inteligencia artificial en los procesos de análisis forense digital para mejorar la identificación, clasificación, correlación e interpretación de evidencia, sin comprometer su integridad, reproducibilidad, explicabilidad y valor probatorio?

%===========================================================
% SECCIÓN 7: OBJETIVOS
%===========================================================
\section{Objetivos}
\subsection{Objetivo General}
Analizar críticamente la integración de técnicas de inteligencia artificial en el ciclo de vida del análisis forense digital para determinar su viabilidad técnica, herramientas disponibles, limitaciones estructurales y condiciones de confiabilidad jurídica en la evaluación de evidencias.

\subsection{Objetivos Específicos}
\begin{enumerate}
    \item \textbf{Identificar} los \textit{frameworks} de referencia actuales, determinando si constituyen estándares industriales o propuestas académicas que combinan la IA con la forensia digital \cite{KhalidZainab2024TauX}.
    \item \textbf{Evaluar} la aplicabilidad de modelos de IA en el triaje de dispositivos móviles y la detección de \textit{deepfakes}, analizando cómo correlacionan artefactos de múltiples fuentes \cite{10.1145/3733102.3733138, VASILARAS2024301737}.
    \item \textbf{Determinar} los controles técnicos, éticos y legales requeridos para que los resultados generados por algoritmos de IA sean trazables, mitigando riesgos de sesgos y sobreajuste, para su admisibilidad en un contexto jurídico \cite{alma991005936804107681, FAHNDRICH2023301617}.
\end{enumerate}

%===========================================================
% SECCIÓN 8: PREGUNTAS DE INVESTIGACIÓN
%===========================================================
\section{Preguntas de investigación}
Para abarcar las exigencias analíticas y metodológicas de este estudio, el desarrollo del artículo da respuesta a las siguientes interrogantes clave:
\begin{itemize}
    \item ¿En qué etapas del proceso forense digital puede incorporarse la IA y qué controles deben implementarse para evitar que altere, contamine o interprete incorrectamente la evidencia original?
    \item ¿Existen \textit{frameworks} de referencia estandarizados que integren IA y forensia digital, y cómo garantizan la trazabilidad de las entradas, transformaciones y salidas del modelo?
    \item ¿Qué riesgos jurídicos introducen los sesgos de entrenamiento, los falsos positivos/negativos y la falta de explicabilidad (modelos de ``caja negra''), y qué función debe conservar el investigador humano en el proceso?
    \item ¿Cómo deben medirse métricas algorítmicas clave (precisión, tasa de error, sensibilidad, \textit{recall}, F1-score y nivel de confianza) para que un perito pueda defender técnicamente el uso de una herramienta ante una autoridad judicial? \cite{WarrenEricM.2025Ceal}
\end{itemize}

%===========================================================
% SECCIÓN 9: METODOLOGÍA
%===========================================================
\section{Metodología de búsqueda y selección bibliográfica}
La presente investigación se fundamentó en una Revisión Sistemática de la Literatura (SLR), estructurada para mapear rigurosamente la convergencia entre la ciberseguridad, la informática forense y la inteligencia artificial explicable.

\subsection{Protocolo de búsqueda y criterios de elegibilidad}
\begin{itemize}
    \item \textbf{Bases de datos consultadas:} Se llevó a cabo un escrutinio exhaustivo en repositorios académicos como ACM Digital Library, IEEE Xplore, y literatura indexada en Scopus, complementado con iRUS. Esta base se trianguló con marcos normativos de entidades como ISO/IEC, NIST, ENFSI y SWGDE.
    \item \textbf{Cadenas de búsqueda utilizadas:} Se emplearon ecuaciones booleanas adaptadas, incluyendo: \textit{``Digital Forensics AND Artificial Intelligence''}, \textit{``Explainable AI (XAI) AND Incident Response''}, \textit{``Deepfake detection AND LLM''}, y \textit{``Mobile forensics AI alignment''}.
    \item \textbf{Periodo cubierto:} La ventana de observación se restringió al quinquenio 2021-2026 para capturar la disrupción de la IA generativa, incluyendo literatura fundacional indispensable previa a esta fecha.
    \item \textbf{Criterios de inclusión:} Artículos científicos revisados por pares, simposios doctorales y reportes técnicos oficiales que aborden arquitecturas de IA, aplicabilidad empírica en forensia y métricas de interpretabilidad.
    \item \textbf{Criterios de exclusión:} Estudios puramente teóricos sin validación, herramientas sin documentación técnica de sus métricas, y artículos que omitieran el análisis de la cadena de custodia.
    \item \textbf{Resultados de la recolección:} Durante la fase de exploración inicial, el \textbf{número de documentos encontrados} fue de 164. Tras aplicar los filtros de exclusión por duplicidad y falta de rigor metodológico, el \textbf{número de documentos seleccionados} de manera definitiva fue de 44 fuentes (comprendiendo 10 \textit{frameworks} arquitectónicos, 20 artículos científicos/empíricos y 14 normativas o estándares técnicos internacionales y legales).
\end{itemize}

%===========================================================
% SECCIÓN 10: ESTADO DEL ARTE
%===========================================================
\section{Estado del arte}
El ecosistema de la informática forense ha experimentado una evolución acelerada, transitando de enfoques heurísticos a la integración de la Inteligencia Artificial a través de múltiples subdominios especializados.

\subsubsection{Triaje Automático y Forensia en Dispositivos Móviles}
El aumento exponencial en el volumen de datos a incautar es el cuello de botella más crítico de la disciplina \cite{QuickDarren2014Ioiv}. Para mitigar este problema, investigadores como Serhal y Le-Khac han implementado algoritmos de aprendizaje automático para perfilar metadatos en dispositivos móviles, logrando priorizar la extracción de evidencia relevante (triaje) y reduciendo drásticamente los tiempos de procesamiento \cite{SerhalCezar2021Mlba}. En el ámbito del Internet de las Cosas (IoT), Kim et al. exploran la adquisición a nivel físico (ej. técnica \textit{chip-off}), demostrando cómo la IA alivia la complejidad de los sistemas operativos cifrados \cite{KimSungbum2022Addf}. No obstante, Vasilaras et al. advierten que la alineación ética y la robustez técnica deben prevalecer, manteniendo al analista humano involucrado para evitar sesgos sistémicos \cite{VASILARAS2024301737}.

\subsubsection{Detección de Medios Sintéticos y Deepfakes}
La proliferación de contenido multimedia manipulado ha requerido contramedidas algorítmicas avanzadas. \textit{Frameworks} multimodelo como DeFaX y arquitecturas basadas en XceptionNet y Vision Transformers no solo detectan \textit{deepfakes} con precisiones superiores al 99\%, sino que integran explicabilidad humana a través de mapas de atención \cite{BharatiNitu2026EddA, Al-ImranMd2025DACF}. Adicionalmente, estudios demuestran el poder de los Modelos de Lenguaje Multimodal (MLLMs) para explicar semánticamente el razonamiento detrás de la detección de imágenes sintetizadas \cite{10.1145/3733102.3733138}, mientras que redes convolucionales especializadas han logrado identificar falsificaciones sutiles de redimensionamiento (\textit{seam carving}) \cite{NevesMiguelJoseDas2025SCFD}. En biometría, el uso de IA en entornos orofaciales 3D ha mejorado la precisión, aunque persisten retos de validación en entornos no controlados \cite{ViegasSofia2026DITf}.

\subsubsection{Reconstrucción Temporal e Inteligencia de Amenazas}
Para la correlación de eventos, la formalización ontológica tradicional \cite{ChabotYoan2014Acfk} y la generación de hipótesis lógicas \cite{CostantiniStefania2015DFEA} han evolucionado hacia arquitecturas neuronales. Trabajos recientes como el de Loumachi et al. demuestran el uso de Generación Aumentada por Recuperación (RAG) para automatizar el análisis de líneas de tiempo de incidentes cibernéticos, correlacionando registros de múltiples fuentes \cite{LoumachiFatmaYasmine2025ACIT}. A nivel de red, el uso de Redes Bayesianas ha permitido modelar la incertidumbre y aislar la actividad adversaria en entornos de tecnología operativa (OT) \cite{MaccaroneLeeT.2025IACi}, mitigando tácticas documentadas en marcos como MITRE ATT\&CK \cite{MundtMichael2022TMoD}.

\subsubsection{IA Explicable (XAI) y Admisibilidad Legal}
Para que los hallazgos de la IA tengan validez judicial, la trazabilidad es no negociable \cite{10.1145/3560107.3560114}. Marcos integrales como XAI-DF proponen flujos de trabajo estandarizados \cite{KhalidZainab2024TauX}, fundamentados en pruebas de concepto que aplican herramientas agnósticas (SHAP y LIME) a la detección de intrusos y discos virtuales \cite{HermosillaPamela2025EAfF, HermosillaPamela2025UoEA, HallStuartW.2022APoC}. Desde la óptica jurídica, Warren et al. argumentan que medir tasas de error mediante métricas estadísticas (como entropía cruzada y F1-Score) es vital para superar filtros de admisibilidad como Daubert o Frye \cite{WarrenEricM.2025Ceal}.

\subsubsection{AI Forensics y Antiforensia}
Un dominio emergente es la \textit{``AI Forensics''}, para investigar a los propios sistemas de IA cuando son vulnerados \cite{alma991005936804107681}. Esto abarca la auditoría de plataformas multiagente para recuperar artefactos forenses \cite{10.1145/3664476.3670908} y la técnica de \textit{``Prompt Scene Investigation''} para rastrear evidencias de \textit{jailbreaks} \cite{10.1145/3774905.3795469}. El avance en esta área resulta crítico para combatir técnicas de antiforensia, tales como la alteración dinámica de marcas de tiempo \cite{10.1145/3712716.3712727}.

%===========================================================
% SECCIÓN 11: FUNDAMENTOS
%===========================================================
\section{Fundamentos de análisis forense digital}
El análisis forense digital es una disciplina técnico-científica fundamentada en un ciclo de vida con fases inmutables: identificación, recolección/adquisición, preservación, análisis y presentación de los datos. La guía NIST SP 800-86 sintetiza este ciclo en cuatro fases: recolección, examen, análisis y reporte \cite{NIST_SP_800_86}. La adquisición de la evidencia exige la aplicación de metodologías estrictamente estandarizadas para asegurar que los artefactos digitales no sufran ningún tipo de alteración (principio de mismidad).

A nivel internacional, la directriz ISO/IEC 27037 \cite{ISO_27037} provee el marco normativo sobre la identificación y preservación inicial de la evidencia, mientras que los laboratorios forenses deben operar bajo los rigurosos estándares de competencia técnica de la norma ISO/IEC 17025 \cite{ISO_17025}. Asimismo, redes científicas como ENFSI y el SWGDE proporcionan los manuales de buenas prácticas para la autenticación de imágenes \cite{ENFSI_BPM_Image_2021}, el análisis de audio \cite{ENFSI_BPM_Audio_2022}, la comparación forense de locutores \cite{ENFSI_BPM_Speaker_2022}, el uso de sistemas de reconocimiento facial \cite{ENFSI_GDL_Facial_2022}, las adquisiciones en la nube \cite{SWGDE_Cloud_2025}, las recolecciones remotas \cite{SWGDE_Remote_2025}, el manejo de dispositivos de almacenamiento dañados \cite{SWGDE_Damaged_2025} y el redimensionamiento de imágenes para procesos judiciales \cite{SWGDE_Resizing_2025}.

La ejecución de una adquisición física requiere el uso de bloqueadores de escritura (hardware o software) y la generación inmediata de firmas criptográficas (mediante funciones \textit{hash} como SHA-256) para certificar la inmutabilidad matemática de la imagen forense. Es imperativo subrayar que \textbf{la Inteligencia Artificial no sustituye ni altera el rigor de la cadena de custodia}; los controles deben garantizar que los algoritmos operen de manera lícita únicamente en la fase de \textit{análisis}, interactuando sobre copias de trabajo (bit a bit) y jamás sobre la evidencia original.

Cuando la IA se emplea para el procesamiento, el Marco de Gestión de Riesgos (NIST AI RMF) exige que estos sistemas sean seguros, transparentes y justos \cite{NIST_AI_RMF_2023}. Los resultados analíticos deben conservar un hilo conductor auditable que traduzca la clasificación del modelo (apoyada en métricas como \textit{precision, recall, F1-score}) en razones de verosimilitud estructuradas \cite{WarrenEricM.2025Ceal}. Este estricto apego a los fundamentos permite a los peritos defender la repetibilidad de los hallazgos y rechazar herramientas de ``caja negra'' que no ofrezcan explicabilidad técnica frente a un juez.

%===========================================================
% SECCIÓN 12: TÉCNICAS DE IA APLICABLES
%===========================================================
\section{Técnicas de inteligencia artificial aplicables}
La implementación de la Inteligencia Artificial en la informática forense exige una categorización taxonómica estricta, diferenciando las capacidades algorítmicas, el manejo dimensional de los datos y los riesgos probatorios inherentes a cada paradigma computacional. A continuación, se analizan críticamente las técnicas aplicables frente a los estándares de admisibilidad pericial y su viabilidad técnica.

\subsection{Aprendizaje Automático Clásico (Machine Learning)}
El aprendizaje automático (ML) tradicional ha demostrado una alta eficacia en el triaje inicial de evidencia y la reducción de dimensionalidad, operando de manera excepcional frente al desafío del \textit{Big Data} forense \cite{QuickDarren2014Ioiv}. La automatización del triaje requiere priorizar artefactos antes de una extracción forense profunda. Los algoritmos supervisados, particularmente aquellos basados en ensambles de árboles de decisión (como \textit{Random Forest}, \textit{Gradient Boosting} y \textit{XGBoost}), se emplean extensivamente en la perfilación de metadatos para dispositivos móviles, clasificación de malware estático y análisis de sistemas de archivos \cite{SerhalCezar2021Mlba}. A diferencia de los enfoques heurísticos rígidos, estos modelos pueden inferir patrones estadísticos complejos a partir del tamaño del archivo, la frecuencia de modificación y las firmas de encabezado.

Sin embargo, desde una perspectiva estrictamente pericial, estos modelos enfrentan desafíos críticos de generalización y equilibrio estadístico. Como evidencian Hermosilla et al., algoritmos predictivos robustos como XGBoost son altamente susceptibles al sobreajuste (\textit{overfitting}) cuando se exponen a conjuntos de datos con severo desbalance de clases, un escenario endémico en la detección de intrusiones de red (NIDS), donde el tráfico anómalo representa una fracción minúscula del flujo total de datos. Este desbalance algorítmico dispara las tasas de falsos positivos y falsos negativos, comprometiendo gravemente la confiabilidad técnica frente a una auditoría judicial o contrainterrogatorio \cite{HermosillaPamela2025UoEA}.

\subsection{Aprendizaje Profundo (Deep Learning)}
El Aprendizaje Profundo (DL) prescinde de la extracción manual de características (\textit{feature engineering}), delegando esta tarea a arquitecturas de redes neuronales multicapa capaces de mapear representaciones latentes jerárquicas.
\begin{itemize}
    \item \textbf{Visión Computacional y Análisis Multimedia:} Las Redes Neuronales Convolucionales (CNN) y las arquitecturas emergentes de \textit{Vision Transformers} (ViT) se aplican para la detección forense de falsificaciones estructurales imperceptibles al ojo humano. Esto resulta vital para combatir técnicas de alteración fotográfica como el \textit{seam carving} (redimensionamiento de imágenes ocultando artefactos) y la proliferación de \textit{deepfakes} \cite{BharatiNitu2026EddA, NevesMiguelJoseDas2025SCFD}. Arquitecturas híbridas avanzadas, como \textit{DeFaX}, han reportado precisiones empíricas superiores al 99\% en la identificación de medios sintéticos generados por Redes Generativas Antagónicas (GAN) y modelos de difusión \cite{Al-ImranMd2025DACF}.
    \item \textbf{Análisis Topológico y de Grafos:} En investigaciones de redes y amenazas avanzadas (APT), el modelado de correlaciones no euclidianas se aborda mediante Redes Neuronales de Grafos (GNN). Estas redes permiten correlacionar el flujo de comunicaciones, el comportamiento temporal de los nodos y los movimientos laterales dentro de ecosistemas corporativos comprometidos, identificando anomalías estructurales de manera automatizada.
\end{itemize}

El riesgo jurídico insoslayable de estas técnicas radica en su naturaleza inherente de ``caja negra''. La extrema dimensionalidad de los tensores subyacentes y la complejidad estocástica de los pesos neuronales ocultan la causalidad lógica de la decisión algorítmica. Esta opacidad estructural impide que el perito testifique con certeza científica sobre \textit{por qué} el algoritmo clasificó un artefacto específico como malicioso o falsificado, invalidando fácticamente su peso probatorio si la inferencia no se acopla a técnicas formales de explicabilidad \cite{VASILARAS2024301737}.

\subsection{Inteligencia Artificial Generativa y Grandes Modelos de Lenguaje (LLMs)}
Los Grandes Modelos de Lenguaje (LLM) y los Modelos de Lenguaje Multimodal (MLLM) representan la frontera computacional actual en la inferencia, correlación e interpretación semántica de evidencia no estructurada. Se ha documentado extensivamente su uso integrado con arquitecturas de Generación Aumentada por Recuperación (RAG) para la reconstrucción automatizada de complejas líneas de tiempo (\textit{timelining}), correlacionando registros de eventos, correos electrónicos y \textit{logs} dispersos en fracciones del tiempo que requeriría un escrutinio manual \cite{LoumachiFatmaYasmine2025ACIT}. De igual manera, los MLLM se están adaptando para extraer y articular explicaciones semánticas detalladas sobre las regiones manipuladas en imágenes sintetizadas \cite{10.1145/3733102.3733138}.

No obstante, su aplicación en laboratorios forenses enfrenta dos barreras operativas y legales severas. En primer lugar, la no reproducibilidad estocástica, derivada de la temperatura del modelo y la configuración de decodificación (\textit{Top-P/Top-K}), riñe directamente con los principios de reproducibilidad científica dictaminados por la norma ISO 17025. En segundo lugar, el riesgo latente de ``alucinaciones'' (invención o extrapolación errónea de datos factuales). Si un modelo genera información ficticia o concatena registros inexistentes al analizar un evento, contamina irreversiblemente la cadena investigativa. Por esta razón, investigadores como Dehing et al. enfatizan de manera categórica que el uso de estas tecnologías debe limitarse exclusivamente a implementaciones locales (\textit{Local LLMs}), operando \textit{offline} (sin conexión a internet) para asegurar la confidencialidad de la evidencia protegida y evitar cualquier posibilidad de exfiltración de telemetría a servidores de terceros \cite{DehingYing2026SRGu}.

\subsection{Inteligencia Artificial Explicable (XAI) y Modelos Probabilísticos Causales}
Para superar de forma sistémica la barrera probatoria y de admisibilidad legal impuesta a los algoritmos de aprendizaje automático, la predicción y el \textit{output} algorítmico deben traducirse, imperativamente, a un lenguaje determinista y auditable por el factor humano. Este requisito ontológico se logra operativamente mediante dos vías metodológicas complementarias:
\begin{enumerate}
    \item \textbf{Modelos Probabilísticos y Causales de Caja Blanca:} El despliegue de arquitecturas transparentes por diseño, destacando el uso de Redes Bayesianas y técnicas de Programación Lógica (como \textit{Answer Set Programming} - ASP). Estos modelos formales cuantifican e integran la incertidumbre matemática en sus inferencias, permitiendo al analista forense trazar y auditar el origen deductivo de un ataque cibernético de forma causal, aislando la actividad adversaria basándose en el teorema de Bayes y la probabilidad condicional pura \cite{MaccaroneLeeT.2025IACi, CostantiniStefania2015DFEA}.
    \item \textbf{Técnicas Post-Hoc de Inteligencia Artificial Explicable (XAI):} La inyección de capas de explicabilidad agnósticas sobre modelos de ``caja negra'' preexistentes. Se destacan métodos de atribución de características como SHAP (fundamentado en la teoría de juegos cooperativos y los valores de Shapley), LIME (modelado sustituto mediante perturbación local de los datos), y Grad-CAM (mapeo topológico de activación visual para CNNs) \cite{HallStuartW.2022APoC, KhalidZainab2024TauX}. La investigación empírica de Hermosilla et al. demuestra contundentemente que la validez científica de estas explicaciones en un tribunal no puede presumirse; debe medirse y documentarse empleando métricas objetivas de robustez analítica, tales como la estabilidad frente a perturbaciones menores (\textit{Stability Score}) y la fidelidad hacia el modelo predictivo subyacente (\textit{Fidelity Score}) \cite{HermosillaPamela2025EAfF}.
\end{enumerate}
Finalmente, para dar cumplimiento irrestricto a estándares jurídicos internacionales en la validación de peritaje científico (como el test de Daubert en Estados Unidos o los criterios de fiabilidad probatoria en jurisdicciones civilistas), Warren et al. argumentan que el desempeño global de estos modelos algorítmicos debe reportarse y presentarse bajo métricas de costo de Razón de Verosimilitud ($C_{llr}$). Esta metodología trasciende las métricas binarias simples y proporciona al juez decisor una medida probabilística y estandarizada del peso empírico y el margen de error pericial \cite{WarrenEricM.2025Ceal}.

%===========================================================
% SECCIÓN 13: FRAMEWORKS Y MODELOS DE REFERENCIA
%===========================================================
\section{Frameworks y modelos de referencia para la integración de IA en el análisis forense digital}
La revisión sistemática permitió identificar diez \textit{frameworks} técnicos (arquitecturas, metodologías, prototipos y modelos) que integran la inteligencia artificial en alguna fase del proceso forense digital, además de tres \textit{frameworks} normativos o de gobernanza que, sin ser específicos de la forensia asistida por IA, delimitan las condiciones en que esos \textit{frameworks} técnicos podrían operar. Para cada uno se estableció quién lo propone, el año de publicación, su campo de aplicación y su categoría: estándar internacional, guía institucional, arquitectura de referencia, modelo de gobernanza, propuesta académica, prototipo experimental o framework implementado en software.

La Tabla~\ref{tab:fw_ident} resume esta identificación. La Tabla~\ref{tab:fw_arq} describe los componentes, entradas, salidas, mecanismos de validación y limitaciones de cada \textit{framework} técnico, y la Tabla~\ref{tab:fw_eval} presenta la matriz comparativa de evaluación. En conjunto, muestran que la mayoría de las soluciones son \textbf{propuestas académicas y prototipos experimentales} diseñados para problemas acotados de un subdominio forense (el triaje de evidencia, la detección de \textit{deepfakes} o la correlación de eventos de red), sin el nivel de validación, despliegue y escrutinio que exigiría un estándar transversal adoptado por la industria \cite{FAHNDRICH2023301617}.

\subsection{Análisis Crítico, Mecanismos Algorítmicos y Nivel de Madurez}
Si bien los marcos rectores emitidos por organismos globales, como el NIST AI RMF \cite{NIST_AI_RMF_2023}, establecen la gobernanza general, la taxonomía de la interpretabilidad y las directrices de mitigación del riesgo algorítmico, a nivel procedimental y táctico-forense las propuestas revisadas se caracterizan por una altísima especialización técnica. A continuación, se desglosa la arquitectura de los diez \textit{frameworks} técnicos identificados:
\begin{itemize}
    \item \textbf{XAI-DF (Khalid, Iqbal y Fung, 2024):} arquitectura académica de referencia que estandariza la aplicación de IA explicable (SHAP y LIME) en investigaciones que involucran volcados de memoria volátil (RAM) y tráfico de red. Como aporte inicial, los autores construyeron una base de datos de 27 volcados de memoria de Windows 7, 10 y 11 con actividad de \textit{malware} simulada, y complementaron la evaluación con el conjunto UNSW-NB15. Su principal innovación procedimental es la integración obligatoria de un \textit{human-in-the-loop}: el investigador actúa como nodo de validación cruzada de las explicaciones y el algoritmo no genera de forma autónoma el informe técnico final \cite{KhalidZainab2024TauX}.
    \item \textbf{DeFaX (Al-Imran et al., 2025):} prototipo experimental que introduce un marco de fusión por atención cruzada (\textit{Cross-Attention Fusion Framework}) entre redes \textit{Swin Transformer}, orientadas al contexto global de la imagen, y modelos \textit{EfficientNet}, orientados a rasgos locales, para detectar medios sintéticos y manipulaciones faciales. Aunque reporta una precisión de 99,8\% en entornos controlados, su transferencia a escenarios judiciales está limitada por su naturaleza de ``caja negra'', que solo se mitiga parcialmente con mapas de calor generados mediante Grad-CAM \cite{Al-ImranMd2025DACF}.
    \item \textbf{GenDFIR (Loumachi, Ghanem y Ferrag, 2025):} propuesta académica para automatizar el análisis de líneas de tiempo de incidentes cibernéticos. Opera en dos etapas: primero preprocesa y estructura los eventos del incidente en una base de conocimiento; luego un agente de Generación Aumentada por Recuperación (RAG) recupera el contexto relevante y un LLM (Llama 3.1 8B en modo \textit{zero-shot}) lo interpreta y enriquece semánticamente. Se evaluó con datos sintéticos en un entorno controlado, con métricas propias y evaluación humana. Aunque reduce el tiempo de análisis, su reproducibilidad pericial es baja por la naturaleza estocástica del modelo generativo subyacente \cite{LoumachiFatmaYasmine2025ACIT}.
    \item \textbf{PSI -- \textit{Prompt Scene Investigation} (Ledjaki et al., 2026):} metodología para auditar y recuperar rastros de sistemas LLM vulnerados mediante inyección de \textit{prompts} (\textit{jailbreaks}) o usados para facilitar actividades ilícitas, un campo emergente denominado \textit{AI Forensics}. A diferencia de los demás \textit{frameworks}, aquí la IA es el objeto investigado y no la herramienta: la guía orienta al investigador para aislar los rastros volátiles (\textit{prompts} de entrada y respuestas del modelo), aplicar \textit{hashing} y preservarlos bajo la doctrina tradicional de cadena de custodia \cite{10.1145/3774905.3795469}.
    \item \textbf{Triaje de archivos en \textit{smartphones} (Serhal y Le-Khac, 2021):} presentado en DFRWS USA 2021, es el \textit{framework} más cercano al alcance de este trabajo en triaje móvil. Propone un flujo de seis etapas (adquisición forense de metadatos, preprocesamiento y normalización, selección de características, ajuste de hiperparámetros, entrenamiento y prueba con validación cruzada de diez particiones, y evaluación del desempeño) para decidir si un archivo extraído de un dispositivo Android debe ser examinado. Su principal fortaleza es que se validó con datos reales: doce dispositivos de casos de terrorismo, con cerca de dos millones de archivos, sobre los que se compararon Naïve Bayes, KNN, SVM, árboles CART, \textit{Random Forest} y redes neuronales. Su debilidad es la contracara de esa fortaleza: los datos no pueden publicarse, de modo que un tercero no puede replicar el experimento \cite{SerhalCezar2021Mlba}.
    \item \textbf{Prueba de concepto de XAI en forensia de discos (Hall, Sakzad y Minagar, 2022):} entrena modelos para predecir el contenido de imágenes y videos y clasificar metadatos del sistema de archivos extraídos de 23 imágenes de disco virtual (VHD) fabricadas, y aplica LIME para explicar las predicciones. Además de clasificar, las explicaciones permiten detectar si el modelo se apoya en características espurias o en sesgos del conjunto de entrenamiento, que es precisamente la verificación que un perito necesitaría para defender un resultado. Su alcance es exploratorio, porque el conjunto de datos es pequeño y artificial \cite{HallStuartW.2022APoC}.
    \item \textbf{Evaluación forense de explicaciones (Hermosilla, Berríos y Allende-Cid, 2025):} más que un sistema de detección, es una metodología para evaluar si las explicaciones de un modelo son útiles en un contexto forense. Aplica SHAP y LIME a dos detectores de intrusiones (XGBoost y TabNet) entrenados con UNSW-NB15 y mide las explicaciones con métricas de fidelidad, estabilidad y similitud de Jaccard. XGBoost alcanzó 97,8\% de exactitud de validación y explicaciones más estables que TabNet. A diferencia de la mayoría de los \textit{frameworks} revisados, cuantifica la calidad de la explicación en lugar de presumirla \cite{HermosillaPamela2025EAfF, HermosillaPamela2025UoEA}.
    \item \textbf{Modelo Bayesiano de Ataque -- BAM (Maccarone et al., 2025):} red bayesiana desarrollada por los laboratorios nacionales de Idaho y Sandia para caracterizar la etapa en que se encuentra un ciberataque contra sistemas de tecnología operativa (OT) a partir de la evidencia observable. Su núcleo es un modelo de proceso de las etapas del comportamiento adversario, y el resto de la red se basa en las tácticas y técnicas de MITRE ATT\&CK para sistemas de control industrial (ICS). Se construyó con datos históricos de 27 ciberataques a sistemas OT documentados en fuentes abiertas y con conocimiento experto, y se demostró con los casos del \textit{ransomware} contra Colonial Pipeline y del ataque al alto horno de Thyssenkrupp. Al ser un modelo de caja blanca, cada probabilidad puede rastrearse hasta un nodo y una evidencia observada; su límite es la escasez de datos abiertos sobre ataques a OT, que obliga a apoyarse en el juicio experto \cite{MaccaroneLeeT.2025IACi}.
    \item \textbf{Programación por Conjuntos de Respuestas -- ASP (Costantini, De Gasperis y Olivieri, 2015):} propone modelar la fase de análisis de evidencia como problemas de lógica computacional resueltos con ASP, de modo que el sistema genere hipótesis investigativas consistentes con los hechos. Esta línea dio origen a la Acción COST CA17124 \textit{DigForASP}, que promueve métodos de IA formales y verificables para el análisis de evidencia. Representa la alternativa simbólica al aprendizaje automático: cada conclusión es deducible y reproducible, pero el modelado exige expertos en lógica y no escala por sí solo a grandes volúmenes de datos \cite{CostantiniStefania2015DFEA}.
    \item \textbf{SADFC -- \textit{Semantic Analysis of Digital Forensic Cases} (Chabot et al., 2014):} enfoque que combina gestión del conocimiento, web semántica y minería de datos para reconstruir líneas de tiempo de incidentes. Su arquitectura gira en torno a una ontología propia: extrae huellas digitales de fuentes heterogéneas, puebla la ontología, aplica operadores formales de reconstrucción y análisis y visualiza la línea de tiempo resultante, buscando la completitud de la reconstrucción. Se incluye como trabajo clásico porque establece un patrón (representación explícita del conocimiento y procedencia de cada evento) que los enfoques recientes basados en LLM, como GenDFIR, todavía no ofrecen \cite{ChabotYoan2014Acfk}.
\end{itemize}

\subsection{Frameworks normativos y de gobernanza}
Los tres \textit{frameworks} de alcance amplio identificados no son procedimientos técnicos de análisis. El \textbf{NIST AI RMF 1.0} (2023) es un modelo de gobernanza voluntario, organizado en cuatro funciones (\textit{Govern}, \textit{Map}, \textit{Measure} y \textit{Manage}), que define las características de una IA confiable, entre ellas que sea válida, explicable e interpretable, pero no prescribe técnicas ni métricas forenses \cite{NIST_AI_RMF_2023}. El \textbf{Reglamento Europeo de IA} (\textit{EU AI Act}, 2024) es un marco regulatorio vinculante que clasifica como de alto riesgo los sistemas usados en la aplicación de la ley y la administración de justicia, y les impone gestión de riesgos, gobernanza de datos, documentación técnica, registro automático de eventos, supervisión humana, exactitud y robustez \cite{EU_AI_Act_2024}. Por último, la norma \textbf{ISO/IEC 27037} es el estándar internacional para identificar, recolectar, adquirir y preservar evidencia digital, pero no contempla el uso de IA \cite{ISO_27037}. Estos \textit{frameworks} establecen \textit{qué} debe garantizarse, pero no \textit{cómo} hacerlo dentro de un caso forense concreto.

%-------------------- TABLA: IDENTIFICACIÓN --------------------
\begin{table*}[!t]
\renewcommand{\arraystretch}{1.25}
\setlength{\tabcolsep}{3pt}
\caption{Identificación y clasificación de los frameworks y modelos de referencia}
\label{tab:fw_ident}
\centering
\scriptsize
\begin{tabular}{|L{2.1cm}|L{2.6cm}|L{2.5cm}|L{3.6cm}|L{3.0cm}|L{2.6cm}|}
\hline
\textbf{Framework} & \textbf{Proponente(s) y año} & \textbf{Categoría} & \textbf{Dominio forense y objetivo} & \textbf{Tipo de IA} & \textbf{Fases cubiertas} \\
\hline
\multicolumn{6}{|l|}{\textit{Frameworks técnicos}} \\
\hline
XAI-DF \cite{KhalidZainab2024TauX} & Khalid, Iqbal y Fung (2024) & Propuesta académica (arquitectura de referencia) & Memoria RAM y red. Unificar el uso de XAI en investigaciones forenses & ML supervisado + XAI agnóstica (LIME, SHAP) & Análisis y presentación \\
\hline
DeFaX \cite{Al-ImranMd2025DACF} & Al-Imran et al. (2025) & Prototipo experimental & Imágenes faciales sintéticas. Detectar \textit{deepfakes} & DL: Swin Transformer + EfficientNet con atención cruzada; Grad-CAM & Análisis \\
\hline
GenDFIR \cite{LoumachiFatmaYasmine2025ACIT} & Loumachi, Ghanem y Ferrag (2025) & Propuesta académica con prototipo & Incidentes cibernéticos (DFIR). Automatizar el análisis de líneas de tiempo & IA generativa: Llama 3.1 8B (\textit{zero-shot}) + agente RAG & Análisis y apoyo al informe \\
\hline
PSI \cite{10.1145/3774905.3795469} & Ledjaki et al. (2026) & Guía metodológica (propuesta académica) & \textit{AI Forensics}. Recuperar evidencia de LLM vulnerados (\textit{jailbreaks}) & La IA es el objeto investigado, no la herramienta & Identificación, adquisición y preservación \\
\hline
Triaje móvil con ML \cite{SerhalCezar2021Mlba} & Serhal y Le-Khac (2021) & Prototipo experimental & \textit{Smartphones} Android. Priorizar archivos de interés a partir de metadatos & ML clásico: NB, KNN, SVM, CART, RF y redes neuronales & Triaje (previo al análisis profundo) \\
\hline
XAI en discos \cite{HallStuartW.2022APoC} & Hall, Sakzad y Minagar (2022) & Prueba de concepto (prototipo) & Discos virtuales (VHD). Clasificar multimedia y metadatos con explicación & ML/DL + LIME & Análisis \\
\hline
Evaluación de XAI en IDS \cite{HermosillaPamela2025EAfF} & Hermosilla, Berríos y Allende-Cid (2025) & Propuesta académica (metodología de evaluación) & Red / detección de intrusiones. Medir la calidad forense de las explicaciones & XGBoost y TabNet + SHAP y LIME & Análisis y validación de resultados \\
\hline
BAM \cite{MaccaroneLeeT.2025IACi} & Maccarone et al., INL y Sandia (2025) & Modelo probabilístico con herramienta en desarrollo & Sistemas OT/ICS. Inferir la etapa de un ataque a partir de observables & Red bayesiana + MITRE ATT\&CK for ICS & Análisis (respuesta a incidentes) \\
\hline
ASP / DigForASP \cite{CostantiniStefania2015DFEA} & Costantini, De Gasperis y Olivieri (2015) & Propuesta académica & Análisis de evidencia. Generar hipótesis investigativas consistentes & IA simbólica: programación lógica (ASP) & Análisis e interpretación \\
\hline
SADFC \cite{ChabotYoan2014Acfk} & Chabot, Bertaux, Nicolle y Kechadi (2014) & Arquitectura de referencia con prototipo & Reconstrucción de eventos. Construir y analizar líneas de tiempo & Ontologías, web semántica y minería de datos & Examen, análisis y presentación \\
\hline
\multicolumn{6}{|l|}{\textit{Frameworks normativos y de gobernanza (referencia)}} \\
\hline
NIST AI RMF 1.0 \cite{NIST_AI_RMF_2023} & NIST (2023) & Modelo de gobernanza (guía institucional voluntaria) & Cualquier sistema de IA. Gestionar riesgos y confiabilidad & No prescribe una técnica & Ciclo de vida del sistema de IA, no del caso forense \\
\hline
EU AI Act \cite{EU_AI_Act_2024} & Parlamento Europeo y Consejo de la UE (2024) & Marco regulatorio vinculante & IA de alto riesgo, incluida la aplicada a la ley y la justicia. Requisitos obligatorios & No prescribe una técnica & Ciclo de vida del sistema de IA \\
\hline
ISO/IEC 27037 \cite{ISO_27037} & ISO/IEC (2012) & Estándar internacional & Evidencia digital. Identificar, recolectar, adquirir y preservar & No contempla IA & Identificación a preservación \\
\hline
\end{tabular}
\end{table*}

%-------------------- TABLA: ARQUITECTURA --------------------
\begin{table*}[!t]
\renewcommand{\arraystretch}{1.25}
\setlength{\tabcolsep}{3pt}
\caption{Componentes, entradas, salidas, validación y limitaciones de los frameworks técnicos}
\label{tab:fw_arq}
\centering
\scriptsize
\begin{tabular}{|L{1.9cm}|L{4.3cm}|L{3.4cm}|L{3.3cm}|L{3.4cm}|}
\hline
\textbf{Framework} & \textbf{Componentes y flujo} & \textbf{Entradas $\rightarrow$ Salidas} & \textbf{Validación experimental} & \textbf{Limitaciones identificadas} \\
\hline
XAI-DF & Base de datos forense $\rightarrow$ extracción de características $\rightarrow$ modelo de ML $\rightarrow$ capa XAI (LIME/SHAP) $\rightarrow$ validación por el investigador & Volcados de RAM y flujos de red $\rightarrow$ clasificación + importancia de características & 27 volcados de memoria (Windows 7, 10 y 11) con \textit{malware} simulado; UNSW-NB15 & Datos simulados; alto costo de calcular valores de Shapley; LIME varía entre ejecuciones \\
\hline
DeFaX & Rama Swin (contexto global) y rama EfficientNet (rasgos locales) $\rightarrow$ fusión por atención cruzada $\rightarrow$ clasificador $\rightarrow$ Grad-CAM & Imagen facial $\rightarrow$ real/sintética + probabilidad + mapa de calor & \textit{Benchmarks} en entorno controlado; 99,8\% de precisión reportada & Generalización a generadores no vistos; sobreajuste al generador de entrenamiento \\
\hline
GenDFIR & Preprocesamiento y estructuración de eventos en una base de conocimiento $\rightarrow$ recuperación de contexto (agente RAG) $\rightarrow$ análisis en lenguaje natural (LLM) & Artefactos y registros del incidente $\rightarrow$ línea de tiempo e informe textual & Datos sintéticos en entorno controlado; métricas propias y evaluación humana & Salidas estocásticas; riesgo de alucinación; sin pruebas con casos reales \\
\hline
PSI & Aislamiento de rastros volátiles (\textit{prompts}, respuestas, \textit{logs}) $\rightarrow$ \textit{hashing} $\rightarrow$ preservación $\rightarrow$ análisis & \textit{Logs} y API del LLM $\rightarrow$ evidencia preservada de las interacciones & Validación procedimental (no evalúa un modelo predictivo) & Requiere acceso a los registros; limitado a LLM \textit{open-source} o con \textit{logs} disponibles \\
\hline
Triaje móvil con ML & Adquisición forense de metadatos $\rightarrow$ preprocesamiento (MinMax) $\rightarrow$ selección de características $\rightarrow$ ajuste de hiperparámetros $\rightarrow$ entrenamiento y prueba $\rightarrow$ evaluación & Metadatos del sistema de archivos $\rightarrow$ archivo ``de interés'' o ``no relevante'' & Validación cruzada de 10 particiones; 12 dispositivos reales, $\approx$2 millones de archivos & Solo Android y un tipo de delito; datos no publicables \\
\hline
XAI en discos & Extracción de multimedia y metadatos de VHD $\rightarrow$ entrenamiento $\rightarrow$ predicción $\rightarrow$ explicaciones LIME & Imágenes, video y metadatos $\rightarrow$ predicción + explicación local & 23 imágenes VHD fabricadas & Conjunto pequeño y artificial; alcance exploratorio \\
\hline
Evaluación de XAI en IDS & Entrenamiento del IDS $\rightarrow$ explicaciones globales y locales (SHAP, LIME) $\rightarrow$ métricas de fidelidad, estabilidad y Jaccard $\rightarrow$ valoración forense & Flujos de red etiquetados $\rightarrow$ alerta + explicación + métricas de calidad & UNSW-NB15; XGBoost con 97,8\% de exactitud de validación & Dataset de laboratorio; el desbalance de clases favorece el sobreajuste \\
\hline
BAM & Modelo de etapas del adversario $\rightarrow$ nodos de tácticas y técnicas ATT\&CK for ICS $\rightarrow$ nodos de observables; probabilidades de datos históricos y juicio experto & Observables del entorno OT $\rightarrow$ probabilidad por etapa de ataque y umbrales de riesgo & 27 ataques a OT de fuentes abiertas; casos Colonial Pipeline y Thyssenkrupp & Pocos datos abiertos; las probabilidades dependen del experto \\
\hline
ASP / DigForASP & Modelado del caso como problema lógico $\rightarrow$ codificación ASP $\rightarrow$ \textit{solver} $\rightarrow$ conjuntos de respuesta (hipótesis) & Hechos extraídos de la evidencia $\rightarrow$ escenarios consistentes con los hechos & Casos de estudio ilustrativos & Requiere expertos en lógica; escalabilidad limitada; no aprende de los datos \\
\hline
SADFC & Extracción de huellas de fuentes heterogéneas $\rightarrow$ poblamiento de la ontología $\rightarrow$ operadores de reconstrucción y análisis $\rightarrow$ visualización & Huellas digitales (registros, navegador, sistema de archivos) $\rightarrow$ línea de tiempo semántica & Prototipo aplicado a escenarios de prueba & Anterior a DL y LLM; mantenimiento de la ontología; escalabilidad \\
\hline
\end{tabular}
\end{table*}

%-------------------- TABLA: MATRIZ DE EVALUACIÓN --------------------
\begin{table*}[!t]
\renewcommand{\arraystretch}{1.25}
\setlength{\tabcolsep}{2.5pt}
\caption{Matriz comparativa de evaluación de los frameworks técnicos (A: alta, M: media, B: baja)}
\label{tab:fw_eval}
\centering
\scriptsize
\begin{tabular}{|L{1.7cm}|L{1.85cm}|L{1.85cm}|L{1.85cm}|L{1.9cm}|L{1.85cm}|L{1.85cm}|L{1.85cm}|L{1.8cm}|}
\hline
\textbf{Framework} & \textbf{Explicabilidad} & \textbf{Intervención humana} & \textbf{Reproducibilidad} & \textbf{Manejo de errores} & \textbf{Trazabilidad} & \textbf{Integridad de la evidencia} & \textbf{Uso judicial} & \textbf{Código o modelo público} \\
\hline
XAI-DF & A: LIME y SHAP & Obligatoria: valida cada explicación & M: LIME exige fijar la semilla & Revisión cruzada humano--explicación & A: registra pesos de características & Opera sobre volcados (copias) & M: auditable, pero validado con datos simulados & Repositorio público reportado \\
\hline
DeFaX & M: Grad-CAM, solo visual & Verifica las regiones señaladas & M: depende de pesos y datos & Falsos negativos ante generadores nuevos & M: mapa de activación & Analiza copias; conserva metadatos & B--M: impugnable por caja negra & Interfaz web (Flask) reportada \\
\hline
GenDFIR & B: justificación textual, no causal & Revisor final obligatorio & B: salidas estocásticas & Alucinación y contaminación de contexto & M: el RAG permite citar eventos recuperados & Riesgo de que el LLM ``complete'' eventos & B: viola la reproducibilidad & Modelo base abierto; código no reportado \\
\hline
PSI & No aplica (no infiere) & El analista dirige el proceso & A: procedimiento documentado & Validación manual estricta & A: \textit{hashing} de \textit{logs} & Bloqueo de escritura y \textit{hashing} & A: respeta la cadena de custodia & Metodología publicada \\
\hline
Triaje móvil con ML & M: importancia de características, sin XAI formal & El perito examina lo priorizado & M: flujo replicable, datos privados & Falsos negativos descartan archivos relevantes & M & Adquisición forense sin alterar la evidencia & M: para priorizar, no como prueba & Datos no públicos (casos reales) \\
\hline
XAI en discos & A: LIME & Interpreta las explicaciones & M: LIME es estocástico & Detecta sesgos y rasgos espurios & M & Opera sobre imágenes de disco & B: exploratorio & No reportado \\
\hline
Evaluación de XAI en IDS & A: explicaciones medidas (fidelidad, estabilidad) & Interpreta las explicaciones & A: dataset público & El desbalance eleva FP y FN & A & No aplica (\textit{benchmark}) & M & Dataset público (UNSW-NB15) \\
\hline
BAM & A: caja blanca & El experto define probabilidades & A: inferencia determinista & Incertidumbre explícita en las probabilidades & A: de cada nodo a su observable & Solo consume observables & M--A: auditable, con supuestos subjetivos & Herramienta documentada; acceso no confirmado \\
\hline
ASP / DigForASP & A: deducción lógica & Modela reglas y valida hipótesis & A: determinista & Errores de modelado, no estadísticos & A & Opera sobre hechos extraídos & M--A: verificable, sin adopción práctica & \textit{Solvers} ASP abiertos \\
\hline
SADFC & A: semántica explícita & Interpreta la línea de tiempo & A & Depende de la calidad de la extracción & A: procedencia de cada evento & Opera sobre huellas extraídas & M & No reportado \\
\hline
\end{tabular}
\end{table*}

\subsection{Matriz comparativa}
La Tabla~\ref{tab:fw_eval} califica cada \textit{framework} técnico con una escala de tres niveles: \textit{alta} cuando el \textit{framework} incorpora el criterio por diseño y lo demuestra; \textit{media} cuando lo cubre parcialmente o depende de la configuración del usuario; y \textit{baja} cuando no lo contempla o lo contradice. Las calificaciones son una valoración de los autores a partir de lo reportado en cada publicación. Junto con las Tablas~\ref{tab:fw_ident} y~\ref{tab:fw_arq}, la matriz cubre las catorce variables de comparación: objetivo, dominio forense, tipo de IA, fases cubiertas, explicabilidad, intervención humana, validación experimental, reproducibilidad, manejo de errores, trazabilidad, protección de la integridad de la evidencia, posibilidad de uso judicial, disponibilidad pública del código o modelo y limitaciones identificadas.

\subsection{¿Existe un framework integral y estandarizado?}
A partir de la matriz, se concluye que \textbf{no existe actualmente un framework integral, estandarizado y ampliamente adoptado} para la forensia digital asistida por IA. Ninguno de los diez \textit{frameworks} técnicos cubre el ciclo forense completo: cada uno atiende una o dos fases y un dominio específico (Tabla~\ref{tab:fw_ident}). Los tres \textit{frameworks} de alcance amplio tampoco llenan ese vacío, porque el NIST AI RMF y el \textit{EU AI Act} fijan requisitos para cualquier sistema de IA sin describir un procedimiento forense, y la ISO/IEC 27037 describe el procedimiento sin contemplar la IA. La comparación revela además cuatro patrones:
\begin{enumerate}
    \item \textbf{Capacidad frente a defendibilidad.} Los modelos con mayor capacidad para procesar datos no estructurados (DeFaX y GenDFIR) obtienen las calificaciones más bajas en explicabilidad y reproducibilidad. Los modelos de caja blanca (BAM, ASP y SADFC) obtienen calificaciones altas en trazabilidad y reproducibilidad, pero dependen del modelado experto y escalan peor. Las técnicas XAI \textit{post-hoc} (XAI-DF, la prueba de concepto de Hall et al. y la evaluación de Hermosilla et al.) ocupan un punto intermedio, y el trabajo de Hermosilla et al. muestra que la calidad de sus explicaciones debe medirse en lugar de presumirse.
    \item \textbf{Validación con datos no reales.} Con excepción del triaje móvil de Serhal y Le-Khac, los \textit{frameworks} se validaron con datos simulados, fabricados o de \textit{benchmark}; y el único validado con casos reales no puede publicar sus datos. Ningún \textit{framework} resuelve la tensión entre realismo y replicabilidad.
    \item \textbf{Métricas de clasificación en lugar de fuerza probatoria.} En los trabajos revisados predominan la exactitud y el F1-score sobre medidas de fuerza probatoria como la razón de verosimilitud o el $C_{llr}$ \cite{WarrenEricM.2025Ceal}.
    \item \textbf{Baja disponibilidad conjunta de código, modelo y datos.} Pocos \textit{frameworks} publican los tres elementos a la vez, lo que impide que una contraparte reproduzca de forma independiente el resultado que se le opone en un proceso judicial.
\end{enumerate}

En consecuencia, los laboratorios periciales siguen dependiendo de la aplicación paralela de normas forenses convencionales, como la ISO/IEC 27037 \cite{ISO_27037}, complementadas con marcos de gestión de riesgos algorítmicos como el NIST AI RMF \cite{NIST_AI_RMF_2023} y, en la Unión Europea, con las obligaciones del \textit{EU AI Act} \cite{EU_AI_Act_2024}. El vacío identificado es un \textit{framework} que articule las tres piezas: un procedimiento técnico por fases, requisitos verificables de trazabilidad y métricas de fuerza probatoria. Esta es la base de la propuesta de mejora presentada en la Sección~\ref{sec:futuro}.

%===========================================================
% DIAGRAMA ARQUITECTÓNICO (TikZ)
%===========================================================
\begin{figure*}[!t]
\centering
\begin{tikzpicture}[
    node distance=1.5cm and 2cm,
    box/.style={rectangle, draw=black, thick, fill=blue!5, text width=3cm, align=center, minimum height=1.5cm, font=\footnotesize\sffamily},
    ai_box/.style={rectangle, draw=red!80!black, thick, fill=red!5, text width=3.2cm, align=center, minimum height=1.5cm, font=\footnotesize\sffamily\bfseries},
    human_box/.style={rectangle, draw=green!60!black, thick, fill=green!5, text width=2.8cm, align=center, minimum height=1.2cm, font=\footnotesize\sffamily},
    arrow/.style={-{Stealth[scale=1.2]}, thick}
]

% Nodos
\node (evidencia) [box] {1. Adquisición\\(Evidencia Original)};
\node (copia) [box, right=of evidencia] {2. Imagen Forense\\(Clon Bit-a-Bit / Hashing)};
\node (ia_model) [ai_box, right=of copia] {3. Inferencia de IA\\(Extracción de Características)};
\node (xai) [ai_box, below=of ia_model] {4. Capa XAI\\(SHAP, Grad-CAM, LIME)};
\node (analista) [human_box, left=of xai] {5. Human-in-the-Loop\\(Revisión y Correlación)};
\node (reporte) [box, left=of analista] {6. Dictamen Pericial\\(Con Razón de Verosimilitud)};

% Conexiones
\draw [arrow] (evidencia) -- (copia);
\draw [arrow] (copia) -- (ia_model);
\draw [arrow] (ia_model) -- (xai);
\draw [arrow] (xai) -- (analista);
\draw [arrow] (analista) -- (reporte);

% Flecha de retroalimentación ortogonal
\draw [arrow, dashed] (analista.north) -- +(0,1.2) -| node[pos=0.25, above, font=\scriptsize\sffamily] {Rechazo/Ajuste de Parámetros} (ia_model.south);

\end{tikzpicture}
\caption{Flujo propuesto de un Pipeline Forense Asistido por IA con validación humana cruzada (Human-in-the-Loop) y capa de explicabilidad (XAI) para preservar la admisibilidad legal.}
\label{fig:pipeline_ia}
\end{figure*}

%===========================================================
% SECCIÓN 14: IMPLICACIONES LEGALES Y ÉTICAS
%===========================================================
\section{Implicaciones legales, éticas y probatorias}
La automatización pericial mediante inteligencia artificial introduce un cambio de paradigma insoslayable en el derecho probatorio y la epistemología judicial. Como se ilustra en la Figura \ref{fig:pipeline_ia}, el flujo de trabajo forense debe incorporar capas de validación humana y explicabilidad algorítmica para evitar que la tecnología socave los pilares constitucionales del debido proceso.

\subsection{El Principio de Contradicción y el Derecho a la Defensa}
Éticamente, el uso de sistemas computacionales opacos traslada la carga de la toma de decisiones críticas a un terreno desprovisto de auditoría independiente. Desde el punto de vista jurídico, fundamentar un dictamen forense en las inferencias de un Gran Modelo de Lenguaje (LLM) o una red neuronal profunda que no pueda explicar causalmente el origen de su deducción vulnera flagrantemente el principio de contradicción de la prueba.

La defensa tiene el derecho inalienable de interrogar la fiabilidad de los métodos empleados para incriminar a un procesado. Si el perito informático desconoce los umbrales de activación neuronal, los pesos específicos de las características evaluadas (\textit{feature importance}) o los sesgos inherentes a la data de entrenamiento, el testimonio experto se desvirtúa \cite{HermosillaPamela2025EAfF}. El analista asume plena responsabilidad civil, penal y disciplinaria sobre su dictamen; por ende, no puede delegar la autoría intelectual del hallazgo a la falibilidad de una máquina \cite{SharmaDevanshu2026LaEI}.

\subsection{Reglamentaciones Nacionales e Internacionales}
La admisibilidad de la evidencia procesada mediante inteligencia artificial debe evaluarse a la luz de los marcos procedimentales locales. En el contexto de Colombia, la Ley 906 de 2004 (Código de Procedimiento Penal) \cite{Ley906_2004} establece en su Artículo 275 los elementos materiales probatorios y la evidencia física, exigiendo en sus artículos subsiguientes la estricta demostración de autenticidad, mismidad e inalterabilidad de la prueba digital. La aplicación de un algoritmo de IA directamente sobre la evidencia física original violaría la cadena de custodia dispuesta en el manual de policía judicial.

En la esfera internacional, normativas emergentes como la Ley de Inteligencia Artificial de la Unión Europea (\textit{EU AI Act}) \cite{EU_AI_Act_2024} catalogan a los sistemas algorítmicos empleados en la administración de justicia y en la aplicación de la ley (\textit{law enforcement}) como sistemas de ``Alto Riesgo''. Esta clasificación impone obligaciones perentorias de trazabilidad, robustez cibernética, supervisión humana ininterrumpida y transparencia, alineándose con las directrices del NIST AI RMF \cite{NIST_AI_RMF_2023}. En este sentido, la IA pericial es lícita únicamente cuando sus resultados estadísticos son empíricamente verificables e interpretables por un operador humano competente e independiente.

%===========================================================
% SECCIÓN 15: LIMITACIONES Y LÍNEAS FUTURAS
%===========================================================
\section{Limitaciones y líneas futuras de investigación}
\label{sec:futuro}
A pesar de los notables avances documentados en el estado del arte, la integración de la IA en los laboratorios forenses enfrenta barreras estructurales que limitan su adopción masiva.

\subsection{Limitaciones Documentadas}
\textbf{1. Disponibilidad de Datasets y Privacidad:} La escasez de conjuntos de datos forenses etiquetados de acceso público representa el obstáculo metodológico más agudo de la disciplina. A diferencia de dominios como la medicina o el procesamiento de imágenes genéricas, la evidencia digital pericial (imágenes de discos, volcados de memoria RAM, capturas de tráfico) está protegida por severas leyes de privacidad, regulaciones de retención de datos (como el GDPR europeo o la Ley 1581 en Colombia \cite{Ley1581_2012}) y el secreto sumarial de investigaciones en curso. Esta limitación obliga a los investigadores académicos a entrenar sus modelos en \textit{datasets} sintéticos, simulados o desactualizados (ej. \textit{Enron Email Corpus}), lo que disminuye drásticamente la precisión del algoritmo al enfrentarse a la entropía y el ruido impredecible de una escena del crimen digital real \cite{QuickDarren2014Ioiv}.

\textbf{2. Obsolescencia Algorítmica y Ataques Adversarios:} La carrera armamentística entre los mecanismos de detección y las técnicas de evasión genera una rápida obsolescencia algorítmica. Por ejemplo, en el dominio del análisis multimedia, los generadores de \textit{deepfakes} impulsados por modelos de difusión (\textit{Diffusion Models}) evolucionan a un ritmo que elude rápidamente los clasificadores forenses entrenados en versiones previas (como las redes GAN tradicionales) \cite{BharatiNitu2026EddA}. Asimismo, los sistemas de IA son vulnerables a ataques adversarios específicos, donde un cibercriminal puede introducir perturbaciones mínimas en el código de un \textit{malware} o en los píxeles de una imagen para forzar deliberadamente un falso negativo en la herramienta de triaje pericial, evadiendo la justicia \cite{MaccaroneLeeT.2025IACi}.

\textbf{3. Cuantificación y Estandarización de la Incertidumbre:} Actualmente, la gran mayoría de las arquitecturas propuestas en la literatura académica reportan su eficacia utilizando métricas de la ciencia de datos computacional pura, tales como la Exactitud (\textit{Accuracy}), el F1-Score, o el Área Bajo la Curva (AUC). Sin embargo, estas métricas son estadísticamente insuficientes para el derecho penal. No logran capturar el nivel de confianza, el margen de error probabilístico ni la incertidumbre epistémica necesarios para superar un escrutinio judicial severo. Un juez no evalúa si un algoritmo tiene un 95\% de \textit{recall}, sino cuál es la probabilidad real de que el hallazgo coincida con la hipótesis acusatoria frente a la hipótesis de la defensa.

\subsection{Propuesta de Mejora y Trabajo Futuro}
Ante el auge de modelos como ChatGPT y Llama, se propone el desarrollo urgente de un marco unificado de gobernanza forense diseñado específicamente para Grandes Modelos de Lenguaje Locales (\textit{Local LLMs}). Este marco debe incorporar por diseño la auditoría forense de interacciones (almacenamiento inmutable de \textit{logs} y \textit{prompts}), garantizando que ninguna evidencia sensible se exfiltre hacia la nube.

Adicionalmente, el trabajo futuro de las ciencias computacionales debe distanciarse de la maximización de la precisión binaria (verdadero/falso) y centrarse en la adopción pericial del Teorema de Bayes. Específicamente, se sugiere estandarizar la métrica de costo de razón de verosimilitud ($C_{llr}$) propuesta por Warren et al. \cite{WarrenEricM.2025Ceal}, estableciéndola como el estándar de oro irrefutable para reportar la tasa de error y la fuerza probatoria de los algoritmos de IA en los tribunales internacionales.

%===========================================================
% SECCIÓN 16: CONCLUSIONES
%===========================================================
\section{Conclusiones}
La integración de la inteligencia artificial en el análisis forense digital no sustituye el rigor procedimental de la disciplina, sino que lo complementa. Si bien las arquitecturas neuronales optimizan el procesamiento de volúmenes masivos de datos y detectan falsificaciones inalcanzables para el ojo humano, su implementación carece de viabilidad jurídica sin una capa robusta de IA Explicable (XAI). Se concluye que ninguna herramienta de IA de ``caja negra'' es idónea para sustentar un hallazgo concluyente en un estrado judicial. La admisibilidad legal dependerá invariablemente de la aplicación del principio \textit{Human-in-the-Loop}, donde la inteligencia artificial opera sobre copias inmaculadas de la evidencia, limitándose a proveer inferencias estadísticas que el perito humano debe corroborar, reproducir y defender técnica y causalmente frente a una autoridad judicial.

\nocite{*}
\printbibliography


\end{document}